% Pista ana-23s-q6 · Anàlisi
% Procedència: PAU Catalunya 2023 · Convocatòria de setembre · Sèrie 2 · Qüestió 6
\documentclass[11pt,a4paper]{article}
\usepackage{pista}
\pistainfo{Catalunya 2023}{Convocatòria de setembre}{Sèrie 2}

\begin{document}

\section*{\textcolor{blauPista}{Qüestió 6}}

\noindent Volem construir una peça metàl·lica que tingui per secció un trapezi isòsceles amb la base superior tres vegades més llarga que la base inferior. Els altres costats del trapezi fan 10 mm, tal com podeu observar en la figura següent:

\noindent\textit{(A la figura: la base superior fa $3x$, la inferior fa $x$ i els dos costats laterals fan 10 mm cadascun.)}

\subsection*{\textit{a)} \normalfont Expresseu l'altura del trapezi en funció de la longitud $x$ de la base inferior. \hspace*{\fill}\textnormal{[0,5~punts]}}

\pista{Com que el trapezi és isòsceles, les dues parts que sobresurten de la base superior respecte de la inferior són iguals: quant fa cadascuna, en funció d'$x$?

\medskip
Dibuixa l'altura $h$ del trapezi. Formes un triangle rectangle amb:
\begin{itemize}
\item catets: l'altura $h$ i una d'aquestes parts;
\item hipotenusa: el costat lateral, de $10\,\text{mm}$.
\end{itemize}
Aplica el teorema de Pitàgores i aïlla $h$.}

\subsection*{\textit{b)} \normalfont Calculeu la longitud de la base inferior del trapezi de manera que l'àrea de la peça sigui màxima i trobeu el valor d'aquesta àrea màxima. \hspace*{\fill}\textnormal{[2~punts]}}

\pista{És un problema d'optimització.
\begin{itemize}
\item \textit{Funció a optimitzar.} L'àrea del trapezi és
\[
A = \frac{(\text{base superior} + \text{base inferior}) \cdot h}{2}
\]
Substitueix-hi les bases $3x$ i $x$, i l'altura $h$ de l'apartat \textit{a)}: obtindràs $A(x)$, en funció d'una sola variable.
\item \textit{Optimitza.} Deriva, resol $A'(x) = 0$ i comprova que el valor trobat és un màxim.
\end{itemize}

\medskip
\textbf{Atenció:} queda't només amb la solució que tingui sentit en el problema, i respon \emph{totes} les parts de la pregunta: la longitud de la base inferior \emph{i} el valor de l'àrea màxima.}

\end{document}
